\documentclass[a4paper,10.5pt]{article}
\usepackage[T1]{fontenc}\usepackage[utf8]{inputenc}
\usepackage[top=2.2cm,bottom=2.2cm,left=2.35cm,right=2.35cm]{geometry}
\usepackage{hyperref}\usepackage{xcolor}\definecolor{accent}{HTML}{1F4E78}\hypersetup{colorlinks=true,urlcolor=accent}
\pagestyle{empty}\setlength{\parindent}{0pt}\setlength{\parskip}{7pt}
\begin{document}
\textbf{Prabhavit Gupta}\\
Venice, Italy \textbar{} +39 351 610 7807 \textbar{} prabhavitg@gmail.com\\[10pt]
Dear Mercor Hiring Team,

I am writing to apply for the AI Safety Red Teamer role. My work over the past three years has been building on top of large language models and watching closely where they get things wrong, which is the instinct this role is asking for.

At UltraStrategy, I built a CV parsing pipeline integrated into the platform: a user uploads a CV, the pipeline extracts the data with Affinda and the Claude API, assesses how relevant the candidate is for specific types of roles, and the result is delivered to the frontend before a standardised Ultra CV is generated. Getting that pipeline reliable meant constantly testing it against messy, ambiguous and adversarial input and checking where the model's output could not be trusted.

At FlowSpace, a workspace marketplace, I built support AI agents that answer customer queries directly from live platform data and escalate to a human whenever their confidence is low. Designing that escalation threshold meant probing the agent with edge cases and unclear questions to find where it would otherwise answer wrongly or overconfidently, which is close to the adversarial evaluation this role needs.

I have strong written communication in English, hands-on experience with prompt design and evaluating model output for failure modes, and I am available immediately for contract work.

I would welcome the opportunity to discuss how I could contribute to Mercor's red teaming work.

Regards,\\
Prabhavit Gupta
\end{document}
